\chapter{Branching bisimilarity}

% =====================================================================
% ESTRUCTURA OBJETIVO DE ESTE CAPÍTULO (ver planificacion-tesis.md)
%   3.1 Definición formal.
%   3.2 Relación con otras equivalencias (weak, observational, WSOE)
%       -> clave para justificar la sustitución del Cap. 4.
%   3.3 Algoritmo O(m log n) de Groote et al. (2019):
%       ideas principales, estructuras (splitters, bloques), invariantes.
%
% Papers de referencia:
%   - Groote, Jansen, Keiren, Wijs (2019), arXiv:1909.10824  -> \cite{groote2019}
%     (copia local en Papers/jansen.pdf)
% =====================================================================

% TODO(intro): párrafo puente desde el Cap. 2. Recordar que en el flujo
% composicional se minimiza entre composiciones usando una noción de
% equivalencia (WSOE, definida en 2.4) y que el costo de esa minimización
% es el cuello de botella. Anunciar que branching bisimilarity es la
% equivalencia candidata a reemplazarla y que este capítulo la define,
% la ubica frente a las equivalencias clásicas y presenta el algoritmo
% eficiente que la calcula.

\section{Definición formal}

% TODO(3.1):
%  - Definir transiciones internas / no observables (acción tau) sobre el
%    LTS ya definido en el Cap. 2; introducir notación si hace falta
%    (reusar \step, \walk, \hop de commands.tex en lugar de inventar).
%  - Definición de branching bisimulation entre dos estados: condición de
%    matching que permite intercalar pasos tau preservando el estado
%    intermedio (la diferencia clave con weak bisimulation).
%  - Branching bisimilarity como la mayor branching bisimulation; mencionar
%    que es relación de equivalencia.
%  - (Opcional) variante divergence-sensitive si se usa en el contexto
%    composicional; aclarar cuál se adopta y por qué.
%  - Ejemplo mínimo: dos LTS branching-bisimilares y un par que NO lo es
%    pero sí es weak-bisimilar (motiva 3.2).

\section{Relación con otras equivalencias}

% TODO(3.2): este es el puente conceptual hacia el Cap. 4.
%  - Ubicar branching bisimilarity en la jerarquía: strong > branching >
%    weak / observational. Dar la intuición de por qué branching es más
%    fina que weak (preserva el potencial de los estados intermedios).
%  - Relación con observational equivalence.
%  - Relación con WSOE (la noción del paper composicional, definida en 2.4):
%    adelantar la afirmación que el Cap. 4 demuestra formalmente y dar la
%    intuición de por qué coinciden en el contexto de control.
%  - Cerrar motivando: si coinciden en lo que importa para síntesis, se
%    puede sustituir WSOE por branching bisimilarity y ganar el algoritmo
%    eficiente de 3.3.

\section{\texorpdfstring{Algoritmo $O(m \log n)$}{Algoritmo O(m log n)} de Groote et al.}

% TODO(3.3): basado en \cite{groote2019} (Papers/jansen.pdf).
%  - Idea general: partition refinement. Partir de una partición inicial
%    de los estados en bloques y refinarla hasta el punto fijo.
%  - Estructuras de datos clave: bloques, partición refinable, splitters,
%    y la maquinaria que da la cota O(m log n) (estrategia "process the
%    smaller half" / constellations-bloques del paper).
%  - Invariantes que mantiene el refinamiento (estabilidad de bloques
%    respecto de los splitters) y por qué garantizan correctitud.
%  - Complejidad: enunciar O(m log n) (m transiciones, n estados) y
%    contrastar con la complejidad de la minimización por WSOE vista en
%    2.4 -> ESTE es el argumento de escalabilidad de toda la tesis.
%  - El pseudocódigo extendido va al Apéndice; acá solo las ideas y
%    estructuras principales (usar el lstdefinelanguage 'pseudocode' de
%    tesis.tex si se incluye algún fragmento).
%  - Nota: la adaptación concreta de este algoritmo al contexto
%    composicional (eventos controlables/no-controlables, estados
%    marcados) y su implementación se tratan en el Cap. 5.
